\documentclass[11pt]{article}

\usepackage[T1]{fontenc}
\usepackage[utf8]{inputenc}
\usepackage[margin=1in]{geometry}
\usepackage{amsmath}
\usepackage{amssymb}
\usepackage{booktabs}
\usepackage{tabularx}
\usepackage{array}
\usepackage{xcolor}
\usepackage{listings}
\usepackage[hidelinks]{hyperref}

% Inline code shorthand: monospaced, with break opportunities after underscores
% so long identifiers/paths do not overflow the text block.
\newcommand{\code}[1]{\texttt{\def\_{\textunderscore\allowbreak}#1}}

% Let TeX stretch the last line a little rather than overflowing the margin.
\setlength{\emergencystretch}{3em}

\lstset{
  basicstyle=\ttfamily\small,
  breaklines=true,
  frame=single,
  framesep=4pt,
  language=Python,
  keywordstyle=\color{blue!60!black},
  commentstyle=\color{gray},
  stringstyle=\color{green!40!black},
  showstringspaces=false,
  columns=fullflexible,
}

\lstdefinestyle{plain}{language={},keywordstyle=\ttfamily,%
  basicstyle=\ttfamily\footnotesize,breaklines=false}
\lstdefinestyle{incar}{language={},keywordstyle=\ttfamily,%
  basicstyle=\ttfamily\small,breaklines=true}

\title{How the DFT Inputs Are Generated\\[2pt]
\large The magnetic-orderings pipeline (\code{atomate2} + \code{jobflow-remote})}
\author{Luca Frey}
\date{\today}

\begin{document}
\maketitle

\section{Scope}

This note describes, end to end, how the VASP input files (\code{INCAR},
\code{KPOINTS}, \code{POSCAR}, \code{POTCAR}) for the magnetism pipeline are
produced from a single input \code{Structure}. Nothing in the repository writes
these files by hand: the pipeline is a thin driver around \code{atomate2}'s
\code{MagneticOrderingsMaker}, which in turn delegates to \code{atomate2}'s VASP
\emph{input-set generators} and, below those, to \code{pymatgen}'s
\code{VaspInputSet}. The sections below follow that delegation chain.

All concrete values quoted here were resolved from the installed packages for the
minimal example in \code{magnetism.py}: rock-salt MnO,
\code{Structure(Lattice.cubic(4.445), ["Mn", "O"], \dots)}.

\section{High-level flow}

The user-facing entry point is \code{magnetism.py}:

\begin{lstlisting}
flow = MagneticOrderingsMaker().make(structure)
submit_flow(flow, worker="magnetism_justus2", project="dft_pipeline2")
\end{lstlisting}

\code{MagneticOrderingsMaker.make()} builds a three-job \code{Flow}:

\begin{enumerate}
  \item \textbf{\code{enumerate orderings}} --- enumerate candidate collinear
        magnetic orderings of the input structure.
  \item \textbf{\code{run orderings}} --- for each ordering, run an (optional)
        relaxation followed by a static calculation. This job replaces itself
        with the actual VASP sub-jobs.
  \item \textbf{\code{postprocess orderings}} --- collect the total energies,
        identify the ground-state ordering, and build a
        \code{MagneticOrderingsDocument}.
\end{enumerate}

Only step~2 generates DFT inputs. The downstream \code{exchange.py} flow
(\code{ExchangeMaker}) consumes the \emph{energies} produced here to fit a
Heisenberg model and run Vampire; it does \emph{not} launch new VASP runs, so it
is out of scope for this note.

\begin{center}
\begin{lstlisting}[style=plain]
Structure
   |
   v  enumerate_magnetic_orderings()          [pymatgen MagneticStructureEnumerator]
ordered structures + origins (fm / afm / ...)
   |
   v  run_ordering_calculations()             [one branch per ordering]
   |
   +--> RelaxMaker.make(struct)   -- RelaxSetGenerator  --> VASP relax
   |        | CONTCAR (prev_dir)
   |        v
   +--> StaticMaker.make(struct)  -- StaticSetGenerator --> VASP static  --> energy
\end{lstlisting}
\end{center}

\section{Step 1: Enumerating magnetic orderings}

\code{enumerate\_magnetic\_orderings} (in \code{atomate2/common/jobs/magnetism.py})
is a thin wrapper around \code{pymatgen}'s \code{MagneticStructureEnumerator}.
With the defaults set by \code{MagneticOrderingsMaker} it is called with

\begin{itemize}
  \item \code{strategies = ("ferromagnetic", "antiferromagnetic")},
  \item \code{automatic = True} (let pymatgen add sensible extra strategies),
  \item \code{truncate\_by\_symmetry = True} (drop very low-symmetry, physically
        implausible orderings),
  \item \code{default\_magmoms = None} (use pymatgen's
        \code{default\_magmoms.yaml}; high-spin initial moments).
\end{itemize}

The enumerator returns a list of ordered \code{Structure} objects together with
an \emph{origin} tag for each (\code{"fm"}, \code{"afm"}, \dots). Crucially, each
returned structure carries a per-site initial magnetic moment (the collinear
$\pm$ pattern that defines the ordering). For MnO this yields \textbf{9
orderings} (1 FM + 8 AFM). These per-site moments are what later become the
\code{MAGMOM} tag --- the magnetic ordering is encoded in the input structure,
not in any INCAR setting written by the pipeline.

\section{Step 2: One relax + one static per ordering}

\code{run\_ordering\_calculations} loops over the enumerated orderings and, for
each, emits up to two VASP jobs (see \code{atomate2/common/jobs/magnetism.py}):

\begin{enumerate}
  \item \code{relax\_maker.make(struct)} --- a structural relaxation of the
        ordered cell. The default \code{relax\_maker} is \code{RelaxMaker()}.
  \item \code{static\_maker.make(relaxed\_structure, prev\_dir=relax.dir\_name)}
        --- a high-quality static energy on the relaxed geometry. The static job
        copies the relaxation outputs (\code{CONTCAR}$\to$\code{POSCAR}, charge
        density, etc.) via \code{copy\_vasp\_outputs} before writing its own
        inputs.
\end{enumerate}

The makers come from \code{MagneticOrderingsMaker}'s defaults:

\begin{lstlisting}
static_maker = StaticMaker(
    input_set_generator=StaticSetGenerator(user_incar_settings={"EDIFF": 1e-7})
)
relax_maker  = RelaxMaker()   # default RelaxSetGenerator
\end{lstlisting}

\noindent So the only project-specific deviation from the \code{atomate2} stock
settings is the tightened \code{EDIFF}~$=10^{-7}$ on the static step (the
energies feed a Heisenberg fit, hence the demand for well-converged total
energies).

Each \code{*Maker} is a \code{BaseVaspMaker}. Its \code{make()} method performs
the same four actions for every job
(\code{atomate2/vasp/jobs/base.py}): (i) optionally copy previous outputs,
(ii) \code{write\_vasp\_input\_set(structure, input\_set\_generator)},
(iii) \code{run\_vasp}, (iv) parse the directory into a \code{TaskDoc}. Step~(ii)
is where the input files are written; it calls
\code{input\_set\_generator.get\_input\_set(structure).write\_input(\dots)}.

\section{Step 3: Assembling one VASP input set}

All of \code{INCAR}, \code{KPOINTS}, \code{POSCAR} and \code{POTCAR} come out of a
single \code{VaspInputGenerator} (subclass of \code{pymatgen}'s
\code{VaspInputSet}). The generator is built from three layers, in increasing
priority.

\subsection{Layer 1 --- the \texttt{atomate2} base config}

The base configuration is defined once in \code{atomate2/vasp/sets/base.py}. It
starts from \code{pymatgen}'s \code{MPScanRelaxSet} config dictionary and then
overrides it back to a \emph{GGA{+}U} set:

\begin{itemize}
  \item \code{METAGGA} removed and \code{GGA = "PS"} --- i.e.\ the
        exchange--correlation functional is \textbf{PBEsol}, not the r$^2$SCAN
        meta-GGA of the parent set. \code{KSPACING} is also removed (a
        \code{KPOINTS} file is used instead).
  \item \code{ALGO = "Fast"}, \code{LREAL = False}.
  \item all \code{LDAU*} keys are inherited from \code{MPRelaxSet}: Hubbard $U$ is
        \textbf{on} (\code{LDAU = True}, \code{LDAUTYPE = 2}) with the
        anion-aware per-element $U$ values (for an oxide, Mn gets
        \code{LDAUU} $= 3.9$\,eV, O gets $0$).
  \item \code{KPOINTS}: automatic mesh by reciprocal density,
        \code{reciprocal\_density = 64} (and \code{= 200} for metals).
  \item \code{POTCAR}: functional \code{PBE\_54} plus a table of per-element
        pseudopotential labels (e.g.\ \code{Mn}$\to$\code{Mn\_pv},
        \code{Fe}$\to$\code{Fe}, \code{V}$\to$\code{V\_sv}); elements not listed
        use the pymatgen default label.
\end{itemize}

\subsection{Layer 2 --- per-job \texttt{incar\_updates}}

Each set generator adds a small dictionary of overrides appropriate to the job
type (\code{atomate2/vasp/sets/core.py}):

\begin{center}
\begin{tabular}{lll}
\toprule
Generator & Purpose & \code{incar\_updates} \\
\midrule
\code{RelaxSetGenerator}  & relaxation & \code{NSW=99, IBRION=2, ISIF=3, LCHARG=False} \\
\code{StaticSetGenerator} & static     & \code{NSW=0, ISMEAR=-5, LORBIT=11, LCHARG=True, LREAL=False} \\
\bottomrule
\end{tabular}
\end{center}

\subsection{Layer 3 --- automatic options and the structure}

\code{pymatgen}'s \code{VaspInputSet.incar} then computes a few values from the
structure itself:

\begin{itemize}
  \item \textbf{\code{MAGMOM}} --- one entry per site, taken directly from the
        ordered structure's per-site moments (Section~3). This is how the FM/AFM
        pattern enters the calculation.
  \item \textbf{\code{ISPIN}} --- set to \code{2} automatically because the
        structure is magnetic (\code{auto\_ispin}).
  \item \textbf{\code{ISMEAR}/\code{SIGMA}} --- \code{auto\_ismear} picks safe
        smearing from the (unknown) band gap; for the static step this is then
        overridden to \code{ISMEAR = -5} by Layer~2.
  \item \textbf{\code{LDAUU/LDAUL/LDAUJ}} --- expanded into per-species lists in
        the site order of the \code{POSCAR}.
\end{itemize}

\subsection{Precedence}

When the final \code{INCAR} is assembled the layers are applied so that, for any
given tag, the winner is (highest first):

\[
\underbrace{\texttt{user\_incar\_settings}}_{\text{e.g. EDIFF}=10^{-7}}
\;>\;
\underbrace{\texttt{incar\_updates}}_{\text{per-job (Layer 2)}}
\;>\;
\underbrace{\texttt{auto\_updates}}_{\text{ISPIN, ISMEAR, \dots}}
\;>\;
\underbrace{\texttt{base config}}_{\text{Layer 1}} .
\]

\code{MAGMOM} and the \code{LDAU*} lists are computed from the structure plus the
config and are not overwritten by the later steps. Finally, because a
\code{KPOINTS} file is present, \code{KSPACING} is stripped from the
\code{INCAR}, and unused tags are pruned.

\subsection{The other three files}

\begin{description}
  \item[\code{POSCAR}] the ordered (and, for the static step, relaxed)
        \code{Structure} passed into \code{make()}.
  \item[\code{KPOINTS}] a $\Gamma$-centred automatic mesh generated from
        \code{reciprocal\_density = 64} for the given lattice (for MnO: a
        $4\times4\times4$ grid).
  \item[\code{POTCAR}] selected by mapping each species through the per-element
        label table and the \code{PBE\_54} functional; the actual potential
        files are read from \code{\$PMG\_VASP\_PSP\_DIR}.
\end{description}

\section{Worked example: static input for an AFM MnO ordering}

For one of the AFM orderings of MnO, \code{StaticSetGenerator(user\_incar\_settings=\{"EDIFF":\,1e-7\})}
produces:

\begin{lstlisting}[style=incar]
ALGO = Fast        EDIFF = 1e-07      ENCUT = 680.0      GGA = Ps
ISMEAR = -5        SIGMA = 0.2        ISPIN = 2          NSW = 0
LDAU = True        LDAUTYPE = 2       LDAUL = 2 0        LDAUU = 3.9 0
LDAUJ = 0 0        LMAXMIX = 4        LORBIT = 11        LREAL = False
LCHARG = True      LAECHG = True      LASPH = True       PREC = Accurate
NELM = 200         MAGMOM = 1*5.0 1*-5.0 2*0.6
\end{lstlisting}

\noindent Reading it back against the layers: \code{EDIFF = 1e-07} is the project
override (Layer~3 user setting); \code{ISMEAR = -5}, \code{NSW = 0},
\code{LORBIT = 11}, \code{LCHARG = True} come from \code{StaticSetGenerator}
(Layer~2); \code{GGA = Ps}, \code{ENCUT = 680}, \code{LDAU = True},
\code{LDAUU = 3.9\,0} come from the \code{atomate2} base config (Layer~1);
\code{ISPIN = 2} and \code{MAGMOM} (one Mn up, one Mn down, two O nearly zero)
are derived from the magnetic ordering. The matching \code{KPOINTS} is a
$\Gamma$-centred $4\times4\times4$ mesh.

\section{Runtime environment}

The numeric INCAR/KPOINTS content above is independent of the cluster, but two
runtime ingredients are supplied by the \code{jobflow-remote} worker
\code{pre\_run} block (\code{.jfremote/dft\_pipeline2.yaml}) and matter for
reproducibility:

\begin{itemize}
  \item \code{PMG\_VASP\_PSP\_DIR} points at a local pseudopotential tree
        (\code{.vasp\_psp\_compat}) and \code{PMG\_DEFAULT\_FUNCTIONAL = PBE},
        which together resolve the \code{POTCAR} files.
  \item \code{ATOMATE2\_VASP\_CMD = "vasp"} (and the $\Gamma$-only variant) plus
        the loaded \code{chem/vasp/6.4.3} module provide the executable that
        \code{run\_vasp} invokes.
\end{itemize}

\section{Summary}

The DFT inputs are fully determined by: (1) the magnetic
\emph{ordering} written into the \code{POSCAR}/\code{MAGMOM} by
\code{MagneticStructureEnumerator}; (2) the \code{atomate2} GGA{+}U base set
(PBEsol, \code{LDAU} on); (3) per-job overrides distinguishing relaxation
(\code{ISIF=3}, \code{NSW=99}) from static (\code{ISMEAR=-5}, \code{NSW=0}); and
(4) the single project tweak \code{EDIFF}~$=10^{-7}$ on the static step. No input
file is templated by hand --- every value is produced by the generator chain
\code{*Maker} $\to$ \code{*SetGenerator} $\to$ \code{pymatgen.VaspInputSet}.

\end{document}
